\documentclass[aspectratio=169,11pt]{beamer}
\usetheme{Madrid}
\usecolortheme{dolphin}
\setbeamertemplate{navigation symbols}{}
\setbeamertemplate{footline}[frame number]
\usepackage[T1]{fontenc}
\usepackage{lmodern}
\usepackage{booktabs}
\usepackage{array}
\usepackage{hyperref}
\providecommand{\repositorysourcebase}{https://github.com/bijux/bijux-pollenomics/blob/main/}
\ExplSyntaxOn
\cs_new_protected:Npn \bijux_source_link:n #1
  {
    \tl_set:Nn \l_tmpa_tl {#1}
    \tl_trim_spaces:N \l_tmpa_tl
    \href{\repositorysourcebase\tl_use:N \l_tmpa_tl}
      {\exp_args:NV \nolinkurl \l_tmpa_tl}
  }
\NewDocumentCommand{\evidence}{m}
  {
    \par\vfill
    {\scriptsize\color{gray}Evidence:\space
      \seq_set_split:Nnn \l_tmpa_seq {;} {#1}
      \seq_map_indexed_inline:Nn \l_tmpa_seq
        {
          \int_compare:nNnT {##1} > {1} {;\allowbreak\space}
          \bijux_source_link:n {##2}
        }
    }
  }
\ExplSyntaxOff
\newcommand{\scope}[1]{\par\medskip{\small\color{blue!65!black}#1}}
\titlegraphic{\small Source-bound research briefing}

\title{Bijux Pollenomics: scientific evidence}
\subtitle{What the methods estimate, what the receipts establish, and what remains open}
\author{Bijux Pollenomics}
\date{2 October 2026}
\begin{document}
\frame{\titlepage}

\begin{frame}{Define the reconstruction terms before interpreting results}
\begin{itemize}
\item \textbf{REVEALS}: Regional Estimates of VEgetation Abundance from Large Sites; estimates regional vegetation fractions from large-site pollen.
\item \textbf{LOVE}: LOcal Vegetation Estimates; uses regional background and small-site pollen to estimate local fractions.
\item The \textbf{Landscape Reconstruction Algorithm (LRA)} combines them. \textbf{PPE} means relative pollen productivity estimate; production, fall speed and transport connect grains to cover.
\end{itemize}
\scope{A standard error (SE) describes the declared model calculation. It is not automatically a confidence interval or a measured-cover error.}
\evidence{reports/science/reveals-theory-paper.md; reports/science/love-theory-paper.md}
\end{frame}

\begin{frame}{Four quantities must not become one accuracy claim}
\small
\begin{tabular}{p{0.23\textwidth}p{0.34\textwidth}p{0.33\textwidth}}
\toprule
Quantity & Denominator and unit & What tests it here? \\
\midrule
Observed pollen & Named grains / eligible grain sum & Withheld-site pollen score \\
Regional REVEALS & Modeled taxon cover / regional target area & Author-code numerical replay \\
Local LOVE & Modeled taxon cover / chosen local area & Declared R/native fixture checks \\
Measured vegetation & Surveyed cover / matched landscape area & No admitted Nordic pair \\
\bottomrule
\end{tabular}
\scope{A low pollen error, a small R/native residual, and a small modeled SE answer different questions.}
\evidence{reports/science/model-use.md}
\end{frame}

\begin{frame}{LRA has a regional step and a local step}
\small
\begin{center}
\fbox{Large-site pollen + PPE + transport + basin}\\[0.5em]
$\downarrow$ REVEALS: regional cover over a declared area\\[0.7em]
\fbox{Regional background + small-site pollen + local ring}\\[0.5em]
$\downarrow$ LOVE: modeled cover within that ring
\end{center}
\begin{itemize}
\item Sugita (2007a) tests regional REVEALS in POLLSCAPE simulations; the accessible abstract stresses sufficiently large, distributed lakes.
\item Sugita (2007b) tests background subtraction and local LOVE in simulations; the relevant local source area remains hard to infer for a past landscape.
\end{itemize}
\scope{The two 2007 simulation papers are not reproduced by the later LandClim author example.}
\evidence{reports/science/reveals-theory-paper.md; reports/science/love-theory-paper.md}
\end{frame}

\begin{frame}{A physical pollen-to-cover model needs more than counts}
\small
\begin{tabular}{p{0.26\textwidth}p{0.61\textwidth}}
\toprule
Input & Scientific role and failure if unspecified \\
\midrule
Age model and window & Assigns a dated depth to an interval; a bin label is not a sample date. \\
Taxon mapping and count sum & Defines which grains enter the response denominator. \\
PPE and fall speed & Correct unequal production and distance-dependent transport. \\
Basin, atmosphere and ring & Set deposition geometry, regional reach and local estimand. \\
Regional cohort & Supplies LOVE background; shared target sites can defeat independence. \\
\bottomrule
\end{tabular}
\scope{Parameter parity under one profile does not qualify those inputs for a different core.}
\evidence{reports/science/landclim-paper.md}
\end{frame}

\begin{frame}{WHH05: one count, age and denominator trace}
\small
\begin{tabular}{p{0.30\textwidth}p{0.57\textwidth}}
\toprule
Source step & Exact retained value \\
\midrule
Physical site & White Horse Hill 2005 lake; grid 512; 150 m radius. \\
Raw depth 16 & 1320 cal BP; 301 grains across 30 named fields. \\
Author preparation & Four named one-grain fields excluded; 297 grains on a 31-target vector; ten taxa positive. \\
Age assignment & 1320 lies in $[1250,1750)$ cal BP; output interval labelled 1500 BP. \\
\bottomrule
\end{tabular}
\scope{The 297-grain denominator belongs to this prepared example row; the 1500 BP label is not a new date.}
\evidence{reports/science/publication-correspondence/author-example-site-replay.md}
\end{frame}

\begin{frame}{WHH05: a modeled site value is not measured cover}
\small
\begin{tabular}{lrr}
\toprule
1500 BP site estimate & Cover fraction & Model SE, fraction \\
\midrule
\emph{Calluna vulgaris} & 0.9248830793 & 0.0468069771 \\
Poaceae & 0.0574176519 & 0.0154854693 \\
\bottomrule
\end{tabular}
\begin{itemize}
\item Author R and native code agree for the 31 WHH05 means and SEs: maxima $4.77\times10^{-15}$ and $8.90\times10^{-16}$ fraction.
\item The nine site vectors enter grid 512's regional calculation; WHH05's site output is not the archived regional vector.
\end{itemize}
\scope{About 92.49\% modeled \emph{Calluna} cover is not 92.49\% of the 297 counted grains.}
\evidence{reports/science/landclim-paper.md}
\end{frame}

\begin{frame}{The LandClim paper publishes regional modeled history}
\small
\begin{itemize}
\item Githumbi et al. (2022) present 31 taxa, stated as 12 plant functional types (PFTs) and three land-cover types (LCTs) over 25 Holocene windows and European one-degree cells.
\item Its 1,128 selected records yield 16,711 aggregated site/window counts, not 16,711 physical sediment depths.
\item The authors' selected-window maps show modeled open-land increase from the mid-Holocene and declining modeled summer-green tree extent. These maps are not a measured-vegetation score or a single-cause attribution.
\end{itemize}
\scope{Original paper result: regional cover estimates with reported model uncertainty; our replay is a separate result.}
\evidence{reports/science/landclim-paper.md}
\end{frame}

\begin{frame}{Author-code parity closes a bounded calculation}
\small
\begin{tabular}{p{0.45\textwidth}p{0.42\textwidth}}
\toprule
Bounded input population & Observed numerical correspondence \\
\midrule
17 example lakes; two grids; 25 bins; 31 taxa & 184/425 possible site-windows supported. \\
Raw-to-prepared source replay & 17 files byte-identical; 871 raw rows; 15,028 prepared cells checked. \\
Independent author-R/native execution & 14,074 supported mean/SE values; max residuals $1.84\times10^{-13}$ and $7.15\times10^{-14}$ fractions. \\
\bottomrule
\end{tabular}
\scope{This checks formatting and arithmetic for the released examples, not the paper's 1,128 historical inputs or vegetation accuracy.}
\evidence{reports/science/method-correspondence.md}
\end{frame}

\begin{frame}{The archive differs on its only exactly aligned window}
\small
\begin{tabular}{lrr}
\toprule
TW2, 100--350 BP; author R vs archive & Grid 512 & Grid 562 \\
\midrule
Largest taxon mean gap, percentage points & 9.19787 & 10.55270 \\
Largest taxon SE gap, percentage points & 3.21615 & 4.25814 \\
31-taxon total variation, cover fraction & 0.092406 & 0.110704 \\
\bottomrule
\end{tabular}
\begin{itemize}
\item Comparison: 31 means and 31 SEs per grid, author fractions converted to archive percentage-cover units; 124 paired values.
\item Other author and archive intervals have different bounds, so they are not treated as identical-time rows.
\end{itemize}
\scope{These are two modeled products. Historical source revisions, preparation and chronology remain unverified; no causal diagnosis follows.}
\evidence{reports/science/publication-correspondence/published-window-provenance.md}
\end{frame}

\begin{frame}{An SE envelope is not a validation interval}
\small
\begin{itemize}
\item At TW2, corresponding mean $\pm$ one-SE envelopes overlap for all 31 taxa in both grids. Unknown dependence and source identity prevent a significance interpretation.
\item Archived ISTS and TBS \emph{mean} sums close over 10,500 grid/window rows. The paper states 12 PFTs, deposit prose 11 and archive header 13 positions.
\item Grouped SEs require their covariance/delta-method route; the archive-producing PFT/LCT code, ISTS meaning and deciduous-\emph{Quercus} mapping are not authenticated.
\end{itemize}
\scope{A matched R residual, a grouped mean identity and a reported SE leave biological interpretation open.}
\evidence{reports/science/publication-correspondence/pft-crosswalk/README.md}
\end{frame}

\begin{frame}{The 2024 LOVE paper and our fixtures have different sites}
\small
\begin{itemize}
\item \textbf{Original 2024 paper.} Eight artificial ponds in a simulated German landscape. LOVEr and LOVEoptim yield local-cover curves; the optimized method is slightly better for inner rings. The Figure 2 layout, generated pollen and random draws are unavailable here.
\item \textbf{Our later DISQOVER fixture.} One local/regional pair and 36 rings. The printed parameters match 30/30 numeric cells by row and 9/10 names exactly. The packaged B\"UK300 example is not the paper's B\"UK200 simulation.
\end{itemize}
\scope{The paper's 100--500 m tendency belongs to its simulation; no Nordic source radius is selected.}
\evidence{reports/science/loveoptim-paper.md}
\end{frame}

\begin{frame}{Finite LOVE output can be invalid vegetation}
\small
\begin{itemize}
\item Constructed interface: 1,000 REVEALS draws $\times$ one 200--300 m ring $\times$ 31 taxa yield 31,000 finite LOVE cells; 8,312 are negative (minimum $-0.46953$).
\item Separate packaged example: only 3/10 first rings satisfy its basin boundary; 16/30 valid taxon/ring cells are negative and six exceed one.
\item The constructed local site also enters its regional cohort, so neither calculation is a held-out ecological comparison.
\end{itemize}
\scope{Negative or greater-than-one fractions remain visible as method diagnostics, never silently normalized into cover.}
\evidence{reports/science/love-theory-paper.md}
\end{frame}

\begin{frame}{English forward pollen: direct baselines are lower here}
\scriptsize
\begin{tabular}{lrrrr}
\toprule
14 paired English sites, 30 taxa & \multicolumn{2}{c}{MAE, fraction} & \multicolumn{2}{c}{Hellinger} \\
Method & Direct & LRA forward & Direct & LRA forward \\
\midrule
Training mean & .020496 & .026273 & .282082 & .360671 \\
Nearest site & .019552 & .021132 & .255767 & .286001 \\
Inverse distance squared (IDW) & .018579 & .021088 & .243673 & .285415 \\
\bottomrule
\end{tabular}
\par\smallskip
\small MAE is mean absolute pollen-fraction error; Hellinger is composition distance from 0 (identical) to 1. Both average equally across sites; lower is better. Direct used all 14 training pollen sites; LRA used a regional anchor and 13 local reconstruction predictors. One of 15 outer folds was refused after three training fits failed. Predictions were frozen before target scoring, but sites were previously explored.
\scope{The score is withheld \emph{pollen} composition, not vegetation reconstruction accuracy or a fair equal-information solver ranking.}
\evidence{reports/science/pollen-prediction.md}
\end{frame}

\begin{frame}{Nordic observed-pollen transfer has support refusals}
\small
\begin{tabular}{lrrr}
\toprule
31-target country exclusion & Physical sites & Scored/attempted bins & Refused bins \\
\midrule
Denmark & 1 & 24/24 & 0 \\
Finland & 28 & 435/435 & 0 \\
Norway & 52 & 1,084/1,178 & 94 \\
Sweden & 3 & 51/54 & 3 \\
\bottomrule
\end{tabular}
\scope{Four prior-exposed country cohorts; these are closed observed-pollen subcompositions, with zero Nordic LRA admissions. Denmark has one target site.}
\evidence{reports/science/source-cohorts.md}
\end{frame}

\begin{frame}{The 16-label and 31-target scores have different denominators}
\small
\begin{itemize}
\item \textbf{31 author taxa.} The Swedish country cohort assigns 221,800/240,966 native grains. Error is measured within this closed author target set.
\item \textbf{16 literal labels.} S\"ambosj\"on covers 224,291/393,822 native grains; Lilla Gloppsj\"on covers 55,638/127,027. Error is measured within each lake's exact-label subcomposition.
\end{itemize}
\scope{A lower MAE on one axis does not establish a better model on the other; excluded and unmapped grains are not zero.}
\evidence{reports/science/pollen-prediction.md}
\end{frame}

\begin{frame}{Independent Nordic vegetation validation is still empty}
\small
\begin{itemize}
\item EMPD2--NFI screen: 16 nearby leads, zero admitted matched measured-vegetation pairs.
\item Three Sm\aa land bogs have published historical maps and pollen anchors, but our retained packet lacks the full distance-weighted map-class fractions needed to score our prediction.
\item A southern Swedish published figure yields 5.11 versus 17.33 percentage-point mean absolute differences on \emph{nine rounded printed labels}; it does not reproduce the original 26-taxon study or score this repository.
\item LandClim archived modeled cover cannot serve as independent measured truth for a related input lake.
\end{itemize}
\scope{A validation pair needs one physical site, matched time, taxa, measured cover and spatial support.}
\evidence{reports/science/vegetation-reference.md}
\end{frame}

\begin{frame}{Resolve source identity, then ask for qualified review}
\small
\begin{enumerate}
\item For grids 512/562 at TW2, request exact historical sequence revisions, ages, count denominator, sites/basins, parameters and producing transformations. Confirm other window definitions and reuse terms.
\item For one proposed Nordic core, bind selected chronology, dated basin geometry, taxon/PPE/fall-speed version and eligible count mass before LRA admission.
\item Ask a palynology/landscape reviewer to assess the physical assumptions; ask a reference owner whether measured cover matches the same site, time, taxa and area.
\item Freeze a new qualified cohort before target access; retain refusals and both direct-pollen and LRA-forward baselines.
\end{enumerate}
\scope{Requests are prepared but unsent. Human scientific reading and independent ecological endorsement remain open.}
\evidence{reports/science/landclim-paper.md; reports/science/model-use.md}
\end{frame}

\appendix
\begin{frame}{Technical appendix: four-country source mass}
\scriptsize
\begin{tabular}{lrrr}
\toprule
31-author-target country & Samples & Assigned/native grains & Scored/attempted bins \\
\midrule
Denmark & 91 & 122,998/129,768 & 24/24 \\
Finland & 1,797 & 1,229,019/1,260,281 & 435/435 \\
Norway & 3,544 & 2,018,295/2,286,020 & 1,084/1,178 \\
Sweden & 228 & 221,800/240,966 & 51/54 \\
\bottomrule
\end{tabular}
\scope{NISP means number of identified specimens (grains here); counts and site-bin outcomes use different units. Assignment is a declared subcomposition, not a complete native flora.}
\evidence{reports/science/source-cohorts.md}
\end{frame}

\begin{frame}{Technical appendix: membership is not prediction support}
\small
\begin{itemize}
\item Neotoma join: 12,388 source sample IDs; 5,660 selected across 84 sites, 6,218 outside selected sites, 233 at selected sites but unselected, and 277 without assigned source country.
\item All 233 same-site unselected samples fail the source pollen screen; six selected Swedish negative-BP samples also fail their source age screen and are not scored.
\item SEAD: 25,109 chronology claims at 853 assigned-country sites. All claim propagation is refused; these are neither pollen observations nor ecological events.
\end{itemize}
\scope{A source row, selected sample, attempted site-bin and admitted prediction are distinct denominators.}
\evidence{reports/science/neotoma-study-membership/README.md; reports/science/source-cohorts.md}
\end{frame}

\begin{frame}{Technical appendix: three further Swedish lakes}
\small
\begin{itemize}
\item Gilltj\"arnen, Lake Avegol and Flinkasj\"on were metadata-only additions outside the original three Swedish targets.
\item 231,115/238,797 native grains map to 31 author targets; 63/63 site-bins yield 189 fixed predictions.
\item Equal-site MAE: IDW 0.017334, nearest 0.017932, training mean 0.020682. Nearest wins at Lake Avegol.
\end{itemize}
\scope{Three prior-exposed lake groups; measured target is pollen, with no LRA or independent vegetation score.}
\evidence{reports/science/sweden-lake-transfer.md}
\end{frame}

\begin{frame}{Technical appendix: Swedish training and source support}
\small
\begin{tabular}{p{0.47\textwidth}p{0.39\textwidth}}
\toprule
Matched 31-target pollen comparison & IDW equal-site MAE \\
\midrule
Two Swedish training sites vs same two + 81 DK/FI/NO; 39 shared bins & 0.032846 vs 0.023226 \\
Same-source policy vs excluding centers within 100 km; 52 shared bins & 0.023151 vs 0.029171 \\
\bottomrule
\end{tabular}
\begin{itemize}
\item Sweden-only refuses 15/54 bins; four-country predicts 54/54.
\item Support-policy study produced 473/486 queries and refused 13; one-bin age shifts and nearest-five cap had smaller, site-dependent effects.
\end{itemize}
\scope{These are three exposed Swedish sites, not a spatial population or calibrated source radius.}
\evidence{reports/science/pollen-prediction.md}
\end{frame}

\begin{frame}{Technical appendix: alternative age support changes eligibility}
\small
\begin{itemize}
\item \textbf{ABB / DALMUT LegacyAge.} Joined 107/107 and 71/72 depths; DALMUT 4.5 cm is unmatched. ABB has eight strict samples in four bins; paired IDW MAE moves .060058 to .059762. DALMUT has two strict samples and 20/22 bins unsupported.
\item \textbf{Sols\o{} LegacyAge.} Only 15/91 comparable samples fall in eight of 24 bins. On 16 labels, nearest MAE .032310 versus .032104 with parent-bin counts; a separate 31-target result is .025673.
\item \textbf{Swedish 1,000 BP source block.} All 54/54 frozen target bins are scored; matched-bin IDW MAE moves .023647 to .024023.
\end{itemize}
\scope{Point ages, marginal ranges and nominal bins are not a joint age posterior or historical LandClim admission.}
\evidence{reports/science/pollen-prediction.md}
\end{frame}

\begin{frame}{Technical appendix: two new Swedish pollen targets}
\small
\begin{itemize}
\item \textbf{S\"ambosj\"on, 16 labels.} All 22/22 bins scored; IDW equal-bin MAE .022991. The 21 interval-contained bins yield .023054 versus original-count .022881. Labels contain 224,291/393,822 native grains; 53/111 samples are interval-contained. The score was inspected before protocol annotation correction.
\item \textbf{Lilla Gloppsj\"on, 16 labels.} All 23/23 bins scored; mean MAE .016272 versus IDW .017295. Nearest Hellinger .228724 versus IDW .237576. Labels contain 55,638/127,027 grains; two samples are age-refused. The target partition was parsed in a prior validation screen.
\end{itemize}
\scope{Two distinct one-site development results; different age support and no untouched confirmation.}
\evidence{reports/science/sambosjon-direct-pollen.md; reports/science/lilla-gloppsjon-direct-pollen.md}
\end{frame}

\begin{frame}{Technical appendix: one-window R lineage}
\small
\begin{itemize}
\item Original 2016 REVEALS function and pinned later Petr Kune\v{s} source ran the same 225 BP input in four site configurations: one lake, unequal pair, six sites, nine-site grid.
\item All 124 paired cover values and 124 SEs agree at retained CSV precision; full-grid bootstrap site indices were identical.
\item Native historical profile agrees with later R on those fixed indices within declared tolerances.
\end{itemize}
\scope{One conditional lineage; distinct R function bodies and one window do not certify corrected uncertainty or ecological validity.}
\evidence{reports/science/method-correspondence.md}
\end{frame}

\begin{frame}{Technical appendix: LOVE numerical correspondence}
\scriptsize
\begin{tabular}{p{0.26\textwidth}p{0.57\textwidth}}
\toprule
Declared source test & Result and limit \\
\midrule
Original/current R and Python & 43/43 core and 23/23 source-domain cases pass; 97 numeric CSVs byte-match native R; 48 original-R shared vectors agree. \\
One fixed ten-taxon LOVEexample & Eight ring/repeat candidate selections agree; 40 finite cover and 40 pollen cells match exactly; max ring-factor/objective gaps $5.11\times10^{-15}$/$4.67\times10^{-15}$. \\
Three shared optimized objectives & Both attain objective gaps $<10^{-4}$ from independent optima; cross-evaluated candidate gap $<10^{-10}$. Search paths and vectors need not match. \\
\bottomrule
\end{tabular}
\scope{Declared arithmetic and objective quality, not the 2024 paper's eight-pond Figure 2 or real vegetation accuracy.}
\evidence{reports/science/method-correspondence.md}
\end{frame}

\begin{frame}{Technical appendix: English LRA routes and refusals}
\scriptsize
\begin{tabular}{p{0.23\textwidth}p{0.29\textwidth}p{0.35\textwidth}}
\toprule
Route & Population & Retained disposition \\
\midrule
Analytical LOVE & 31 taxa; 15 English folds & 15/15 invalid compositions; all 54 local ring rows include negative cover. \\
Differential evolution & 30 taxa; one pair, two requests & Both refused convergence; no held-out score. \\
Constrained SLSQP & 30 taxa; 15 folds, 195 training items & 192/195 fits complete; BECKHAM outer fold refused; 14 folds yield 42 scored predictions. \\
\bottomrule
\end{tabular}
\scope{SLSQP is sequential least-squares programming. The selected 1250--1750 BP profile excludes 159 Filipendula grains from 20,732 source NISP; routes do not form a paired solver ranking.}
\evidence{reports/science/english-lra-pollen-prediction.md}
\end{frame}

\begin{frame}{Technical appendix: target basin radius is a forward sensitivity}
\small
\begin{tabular}{lrrr}
\toprule
English target radius factor & 80\% & 100\% & 120\% \\
\midrule
IDW-vegetation forward pollen MAE & .021032 & .021088 & .021457 \\
\bottomrule
\end{tabular}
\begin{itemize}
\item All 42 baseline LRA pollen fractions replay exactly before scoring; 126 radius projections use the same 14 target sites.
\item Vegetation, regional cover and PPE are fixed, so only the forward consumer's assumed lake radius changes.
\end{itemize}
\scope{The shift exposes geometry sensitivity; it does not calibrate a physical radius or admit a Nordic basin.}
\evidence{reports/science/english-lra-pollen-prediction.md}
\end{frame}

\begin{frame}{Technical appendix: tuning cannot be promoted across axes}
\scriptsize
\begin{tabular}{p{0.37\textwidth}p{0.42\textwidth}}
\toprule
Development search & Matched result \\
\midrule
17-site radial basis function (RBF) vs fixed IDW, 31 taxa & 174 common bins: MAE .017774 vs .018948; RBF covers 174/184 attempted vs IDW 178/184. \\
Country-excluded 16-label inner-fold RBF selection & Finland selected .034871 vs IDW .034375; Norway .033942 vs IDW .027012; Denmark selected IDW .033022. \\
Two-case synthetic PPE/transport sweep & Winner reflects a transport-law mismatch, not a real-core PPE recommendation. \\
\bottomrule
\end{tabular}
\scope{Coverage, response axis and target exposure govern model use; no universal LRA parameter follows.}
\evidence{reports/science/model-use.md}
\end{frame}

\begin{frame}{Technical appendix: context cannot fill the reference gap}
\small
\begin{itemize}
\item HYDE modeled cropland: 84 country/year/scenario grid sums match retained source totals within $8.51\times10^{-11}$ km$^2$; rights/version and historical spatial semantics remain open.
\item LandClim Nordic map: 1,875 normalized published-cell features are modeled estimates, not independent field vegetation.
\item Named core gaps: Abborrtj\"arnen needs selected chronology/basin/denominator/PPE; Sols\o{}, Dalmutladdo and Laihalampi have source links but lack complete time-matched LRA input closure.
\item The next ecological test still requires a separately governed measured-cover pair and qualified review.
\end{itemize}
\scope{A clean map export or historical context layer does not add an admitted vegetation truth label.}
\evidence{reports/science/source-cohorts.md; reports/science/model-use.md}
\end{frame}
\begin{frame}{Technical appendix: country inputs still missing}
\small
\begin{tabular}{p{0.14\textwidth}p{0.76\textwidth}}
\toprule
Country & Named-core acceptance still needed \\
\midrule
Denmark & Sols\o{}: original 79 input rows, selected chronology and dated basin. \\
Finland & Laihalampi: selected age vector and historical basin for the 150-sample lead. \\
Norway & Dalmutladdo: historical basin and qualified count/parameter/use closure despite matched source controls. \\
Sweden & ABB/ABBO, BUS and BOKE: one jointly qualified core/window packet, including selected ages and basin history. \\
\bottomrule
\end{tabular}
\scope{The 31-target country studies score 1,594/1,691 bins with 97 refusals. They admit zero LRA samples; no independent Nordic measured-cover benchmark or untouched confirmation follows.}
\evidence{reports/science/country-evidence-status.md}
\end{frame}
\end{document}
